\subsection{弱拓扑的乘积}

命题 8. --- 设 \((\mathbf{F}_{\iota}, \mathbf{G}_{\iota})_{\iota \in \mathbf{I}}\) 为一族处于对偶的空间对。设 \(\mathbf{F} = \prod_{\iota \in \mathbf{I}} \mathbf{F}_{\iota}\) 为诸 \(\mathbf{F}_{\iota}\) 的乘积空间，\(\mathbf{G} = \bigoplus_{\iota \in \mathbf{I}} \mathbf{G}_{\iota}\) 为诸 \(\mathbf{G}_{\iota}\) 的直和。若对所有 \(x = (x_{\iota}) \in \mathbf{F}\) 与所有 \(y = (y_{\iota}) \in \mathbf{G}\)，我们写 \(\langle x, y \rangle = \sum_{\iota \in \mathbf{I}} \langle x_{\iota}, y_{\iota} \rangle\)（这是一个只有有限多个非零项的和），则拓扑 \(\sigma(\mathbf{F}, \mathbf{G})\)（相对于双线性型 \((x, y) \mapsto \langle x, y \rangle\)）是诸拓扑 \(\sigma(\mathbf{F}_{\iota}, \mathbf{G}_{\iota})\) 的乘积。

因为，给定一个拓扑 \(\mathcal{T}\) 于 \(\mathbf{F}\) 上；为使对所有 \(y \in \mathbf{G}\) 线性型 \(x \mapsto \langle x, y \rangle\) 都关于 \(\mathcal{T}\) 连续，按 \(\langle x, y \rangle\) 的定义，其充分必要条件是：每个映射 \(x \mapsto \langle \mathrm{pr}_i x, y_i \rangle\) 都关于 \(\mathcal{T}\) 连续，其中 \(i\) 在 I 中任意，而 \(y_i\) 在 \(\mathbf{G}_i\) 中任意；而这意味着每个映射 \(\mathrm{pr}_i\)（F 到 \(\mathbf{F}_i\) 中）关于 \(\mathcal{T}\) 与 \(\sigma(\mathbf{F}_i, \mathbf{G}_i)\) 都连续（I, p.~10, 推论 1）；证毕。

注. --- F 与 G 之间的对偶在 F 中（相应地，在 G 中）分离，当且仅当对所有 \(\iota\in I\)，\(F_{\iota}\) 与 \(G_{\iota}\) 之间的对偶在 \(F_{\iota}\) 中（相应地，在 \(G_{\iota}\) 中）分离。若 F 与 G 之间的对偶在 F 中（相应地，在 G 中）分离，则在 F 中（相应地，在 G 中），与某个 \(G_{\iota}\)（相应地，\(F_{\iota}\)）正交的子空间——它被典范地等同于 G（相应地，F）的一个子空间——是诸 \(F_{\kappa}\)（\(\kappa\neq\iota\)）的乘积所成的子空间（相应地，诸 \(G_{\kappa}\)（\(\kappa\neq\iota\)）的直和）。

推论 1. --- 设 F 与 G 为两个处于分离对偶的向量空间。若空间 F（赋 \(\sigma(F, G)\)）是两个子空间 M、N 的拓扑直和，则空间 G（赋 \(\sigma(G, F)\)）是分别与 M 和 N 正交的子空间 M°、N° 的拓扑直和。

因为，设 \(p: \mathbf{F} \to \mathbf{M}\)、\(q: \mathbf{F} \to \mathbf{N}\) 为对应于 \(\mathbf{F}\) 分解为 \(\mathbf{M}\) 与 \(\mathbf{N}\) 之直和的投影；在这些条件下，映射 \((p, q): \mathbf{F} \to \mathbf{M} \times \mathbf{N}\) 是拓扑同构。若 \(\mathbf{M}_1 = \mathbf{G} / \mathbf{M}^\circ\)、\(\mathbf{N}_1 = \mathbf{G} / \mathbf{N}^\circ\)，则 \(\mathbf{M}\) 与 \(\mathbf{N}\) 上的拓扑（由 \(\mathbf{F}\) 的拓扑诱导）分别与 \(\sigma(\mathbf{M}, \mathbf{M}_1)\)、\(\sigma(\mathbf{N}, \mathbf{N}_1)\) 一致（II, p.~48, 命题 7）。映射 \(^t(p, q): \mathbf{M}_1 \times \mathbf{N}_1 \to \mathbf{G}\)，当给 \(\mathbf{M}_1, \mathbf{N}_1\) 与 \(\mathbf{G}\) 分别赋以拓扑 \(\sigma(\mathbf{M}_1, \mathbf{M})\)、\(\sigma(\mathbf{N}_1, \mathbf{N})\) 与 \(\sigma(\mathbf{G}, \mathbf{F})\) 时，由命题 8 是一个拓扑同构。在这个映射下，\(\mathbf{M}_1\)（相应地，\(\mathbf{N}_1\)）在 \(\mathbf{G}\) 中的像是子空间 \(\mathbf{N}^\circ\)（相应地，\(\mathbf{M}^\circ\)），而拓扑 \(\sigma(\mathbf{M}_1, \mathbf{M})\)（相应地，\(\sigma(\mathbf{N}_1, \mathbf{N})\)）的像是 \(\mathbf{N}^\circ\)（相应地，\(\mathbf{M}^\circ\)）上由 \(\sigma(\mathbf{G}, \mathbf{F})\) 诱导的拓扑，由此推论得证。

推论 2. --- 设 \((e_{\iota})_{\iota \in \mathbb{J}}\) 为向量空间 F 的一组基，F* 是 F 的对偶，并设 \(u: \mathbf{R}^{(I)} \to \mathbf{F}\) 为由这组基定义的（代数）同构。则转置映射 \({}^t u: \mathbf{F}^* \to \mathbf{R}^I\) 当 F* 赋拓扑 \(\sigma(\mathbf{F}^*, \mathbf{F})\)、\(\mathbf{R}^I\) 赋乘积拓扑时是拓扑同构。

我们知道（A, II, § 2.6, 命题 10），\({}^t u\) 是双射，并且若对某个 \(x^* \in \mathbf{F}^*\) 令 \(\langle e_i, x^* \rangle = \xi_i^*\)（对所有 \(i \in I\)），则像 \({}^t u(x^*)\) 是向量 \((\xi_i^*)\)，它属于 \(\mathbf{R}^I\)，从而对所有 \(x = \sum_{i} \xi_i e_i\)（在 \(F\) 中），有 \(\langle x, x^* \rangle = \sum_{i \in I} \xi_i \xi_i^*\)。推论随即由这个公式与命题 8 得出。

